% TOP_PROBLEM: {"id": 20001084, "title": "Sliding progression-free windows via gap words", "category_id": 2, "category": {"id": 2, "name": "combinatorics", "display_name": "Combinatorics", "description": "Counting problems, graph theory, discrete structures.", "slug": "combinatorics", "order_index": 2, "created_at": "2026-07-31T15:26:25.671Z"}, "status": "partially_solved", "problem_number": "AIM-COMBINATORICS-0209", "set_id": 14, "statusQualification": "The source leaves density unspecified; upper asymptotic density is defined, and the natural-density variant is explicitly retained."}
% FIRST_POSTED: 2026-10-01
% ENTRY_KIND: statement_only
% UnsolvedMath ID 20001084; source catalogue code AIM-COMBINATORICS-0209
% !TeX program = lualatex
% Definition and sources only; this file is not an LLM solution attempt.
\documentclass[11pt,a4paper]{article}
\usepackage[margin=25mm]{geometry}
\usepackage{fontspec}
\setmainfont{lmroman10-regular.otf}[BoldFont=lmroman10-bold.otf,
 ItalicFont=lmroman10-italic.otf,BoldItalicFont=lmroman10-bolditalic.otf]
\usepackage{amsmath,amssymb,mathtools}
\usepackage{unicode-math}
\setmathfont{latinmodern-math.otf}
\AtBeginDocument{\let\setminus\smallsetminus}
\usepackage{xurl,hyperref}
\hypersetup{colorlinks=true,urlcolor=blue}
\setcounter{secnumdepth}{0}
\setlength{\parindent}{0pt}
\setlength{\parskip}{5pt}
\setlength{\emergencystretch}{3em}
\begin{document}
\section*{Sliding progression-free windows via gap words}\label{20001084}
{\small UnsolvedMath ID 20001084 · AIM-COMBINATORICS-0209 · Combinatorics · AIM Workshop Problem Lists}
\subsection{Definitions and mathematical statement}
Fix an integer $s\ge3$. Let $A=\{a_1<a_2<\cdots\}$ be an infinite
increasing sequence of nonnegative integers. Require that every block
of $s$ consecutive terms, $\{a_{i+1},\ldots,a_{i+s}\}$ for $i\ge0$,
contain no nontrivial three-term arithmetic progression.
A nontrivial progression consists of distinct elements $x<y<z$
satisfying $x+z=2y$. The block is consecutive in the ordered list of
members of $A$, rather than a block of $s$ consecutive ambient integers.
Equivalently, a progression $a_i,a_j,a_k$ with $i<j<k$ is allowed
only if $k-i\ge s$.

To specify density for sequences without a natural density, set
\[
 \overline d(A)=\limsup_{N\to\infty}
                \frac{|A\cap\{0,1,\ldots,N-1\}|}{N},\qquad
 \Delta_s=\sup_A\overline d(A),
\]
where the supremum runs over sequences satisfying the block condition.
The problem is to determine $\Delta_s$ and describe the extremal sets
attaining it. If density is required to exist as a limit, the parallel
question uses that natural density instead; this distinction is stated
explicitly because the source does not specify a convention.
The given examples are the residue classes $0,1,3,4$ modulo $6$ for
$s=4$, of density $2/3$, and $0,1,3,4,9,10,12,13$ modulo $18$ for
$s=8$, of density $4/9$. These provide benchmarks, not an assertion
that all extremizers are periodic.
\subsection{Short English statement}
How dense can an integer sequence be if each block of s consecutive sequence elements avoids three-term progressions, and what do extremizers look like?
\subsection{Sources}
\begin{itemize}
\item[S1] ULAM.AI and the UnsolvedMath Contributors, supplied problems.json, record 20001084 (AIM-COMBINATORICS-0209), AIM Workshop Problem Lists. Catalogue identity and source transcription; CC BY 4.0.
\url{https://huggingface.co/datasets/ulamai/UnsolvedMath}.
\item[S2] American Institute of Mathematics, Recent trends in additive combinatorics, item 1.3. Original workshop question underlying AIM-COMBINATORICS-0209.
\url{https://aimath.org/WWN/additivecomb/additivecomb.pdf}.
\end{itemize}
\end{document}
